\chapter{Machine Learning Methodology}\label{ch:method}

\section{Problem formulation and metrics}
Direct forecasting learns $f_h(x_{\le t})\approx y_{t+h}$. Because a measured level is very persistent over minutes, the project also evaluates delta forecasting,
\begin{equation}\label{eq:delta}
\Delta y_{t,h}=y_{t+h}-y_t,\qquad \hat y_{t+h}=y_t+\widehat{\Delta y}_{t,h},
\end{equation}
which concentrates the model on future \emph{movement} and is valid only when the laboratory value is fresh. Target-free virtual sensing learns $g(x_{\le t})\approx y_t$ without any current, past or future target feature.

Errors are reported with MAE, RMSE and $R^2$. Forecasts are compared with the operational persistence baseline $\hat y_{t+h}=y_t$ through the relative RMSE improvement
\begin{equation}\label{eq:improvement}
I_h=100\Bigl(1-\frac{\mathrm{RMSE}_{\text{model}}}{\mathrm{RMSE}_{\text{persistence}}}\Bigr),
\end{equation}
which is negative when the model is worse than persistence. RMSE weighs large errors more and highlights unstable periods; MAE is directly readable as an average deviation.

\section{Chronological validation}
The baseline split is 70\%/15\%/15\% in time order (Table~\ref{tab:splits}); shuffling is excluded because adjacent minutes are nearly identical states and would leak into the test set \citep{bergmeir2018forecasting}. Lag selection, clipping bounds, correlation filtering, hyperparameters and blend weights were learned without test data. Stability research then used four expanding blocks of 1,657 validation rows each (Table~\ref{tab:folds}); within each fold, imputation, scaling, clustering and local experts are fitted on the preceding rows only, giving a distribution of errors rather than a single number.

\begin{figure}[!htbp]\centering
\includegraphics[width=0.95\textwidth]{feature_preparation_artifacts/figures/split_timeline.png}
\caption{Chronological train/validation/test partition of the January series.}\label{fig:split}
\end{figure}

Figure~\ref{fig:split} shows the partition; the validation and test periods always follow the training period in time.

\begin{table}[!htbp]\centering\small\zebra
\caption{Expanding chronological folds used for stability analysis (UTC).}\label{tab:folds}
\begin{tabularx}{\textwidth}{L{0.08\textwidth}R{0.15\textwidth}Y}\toprule
\thead\hd{Fold} & \hd{Train rows} & \hd{Validation block}\\\midrule
1 & 30,940 & 2026-01-22 17:18 to 2026-01-23 21:08\\
2 & 32,597 & 2026-01-23 21:09 to 2026-01-25 00:49\\
3 & 34,254 & 2026-01-25 00:50 to 2026-01-26 04:30\\
4 & 35,911 & 2026-01-26 04:31 to 2026-01-27 08:58\\\bottomrule
\end{tabularx}
\end{table}

\section{Models and tuning}
Ridge regression, $\hat\beta=\arg\min_\beta\lVert y-X\beta\rVert_2^2+\alpha\lVert\beta\rVert_2^2$ on standardised features, is the regularised linear reference for correlated engineered features \citep{hoerl1970ridge}. Extra Trees \citep{geurts2006extratrees} and histogram gradient boosting test nonlinear interactions; the target-free study adds PLS \citep{wold2001pls}, KNN, XGBoost, LightGBM \citep{ke2017lightgbm} and validation-weighted blends. Each family received a small, comparable search budget (two candidates per model, Table~\ref{tab:hyper}); the validation split selects the candidate.

\begin{table}[!htbp]\centering\small\zebra
\caption{Search space of the target-free benchmark (Notebook 05).}\label{tab:hyper}
\begin{tabularx}{\textwidth}{L{0.20\textwidth}YY}\toprule
\thead\hd{Model} & \hd{Candidate 1} & \hd{Candidate 2}\\\midrule
Ridge & $\alpha=1$ & $\alpha=10$\\
PLS & 5 components & 10 components\\
KNN & $k=25$, distance weights & $k=75$, distance weights\\
Extra Trees & 60 trees, leaf 2, features 0.8 & 60 trees, leaf 5, features 0.7\\
Hist. gradient boosting & 80 iter., rate 0.05, 15 leaves & 80 iter., rate 0.08, 31 leaves\\
XGBoost & 150 trees, depth 4, rate 0.05 & 150 trees, depth 6, rate 0.05\\
LightGBM & 150 trees, 15 leaves, rate 0.05 & 150 trees, 31 leaves, rate 0.05\\\bottomrule
\end{tabularx}
\end{table}

The robustness notebooks compare global blends, KMeans-partitioned local Ridge experts (a cluster is assigned to each sample by nearest centroid in standardised feature space), a global/local hybrid and just-in-time learning (JITL) with PCA or PLS neighbourhoods \citep{cheng2004jitl}. KMeans clusters are statistical groups of feature space, not confirmed operating regimes. The stability-first selection rule ranks candidates by mean RMSE, RMSE standard deviation across folds and worst-fold RMSE.

\section{Uncertainty and drift flag}
For the virtual sensor, an empirical 90\% interval is $\hat y_t\pm q_{0.9}$, where $q_{0.9}$ is the 90\textsuperscript{th} percentile of absolute validation residuals of the individual Ridge model (0.507). A drift flag marks samples lying outside the training envelope of the inputs. Both are advisory diagnostics, not guarantees.
